\documentclass[10pt,oneside]{book}
%%%% Page Info + Commands %%%%%{

%packages
\usepackage{pgfplots}
\usepackage{mathrsfs}
\pgfplotsset{compat=1.18}
\usepackage{amsmath}
\usepackage{amsfonts}
\usepackage{charter}
\usepackage{amsthm,letltxmacro}
\usepackage{amssymb}
\usepackage{fancyhdr}
\usepackage{float}
\usepackage[most]{tcolorbox}
\usepackage{enumerate}
\usepackage{enumitem}
\usepackage[dvipsnames]{xcolor}
\usepackage{changepage}
\usepackage{mdframed}
\usepackage{graphicx}

% SMALL PAGE COMMAND
\newcommand{\smallpage}{  \setlength \oddsidemargin{.25in}
            \setlength \textwidth{5.5in}
            \setlength \topmargin{0in}
            \setlength \textheight{9in}}


% Commands for sets of numbers
\newcommand{\Z}{\mathbb{Z}}\newcommand{\R}{\mathbb{R}}\newcommand{\N}{\mathbb{N}}\newcommand{\Q}{\mathbb{Q}}\newcommand{\C}{\mathbb{C}}
\newcommand{\real}[1]{\text{Re}\left(#1\right)}
\newcommand{\im}[1]{\text{Im}\left(#1\right)}

% gordie spacing and boxing commands
\newcommand{\gspc}{\vspace{0.13cm}} % 0.5 space
\newcommand{\gline}[0]{\gspc\\} % 1.5 space
\newcommand{\gbline}{\gspc\gspc\\} % double space
\newcommand{\gbox}[1]{\fbox{\parbox{\linewidth}{#1}}} % multiline box command
\newcommand{\gmod}[1]{\,(\text{mod}\,#1)}


\newcommand{\gimage}[2]{
\begin{figure}[H]
    \centering
    \includegraphics[width=#2]{#1}
\end{figure}
}

\newcommand{\centerit}[1]{{\begin{center} #1 \end{center}}}
\newcommand{\ul}[1]{\underline{#1}}

\newcommand{\prt}[2]{\frac{\partial#1}{\partial#2}}

%%%%%%%% command for graphics %%%%%%%%%%%%%

%}

% COLOR BOX
\newmdenv[backgroundcolor=gray!8,
  linewidth=0pt, innerleftmargin=0pt, innerrightmargin=0pt,
  skipabove=\baselineskip, skipbelow=\baselineskip
]{coloredadjust}

% PROOF
\LetLtxMacro\oldproof\proof
\let\endoldproof\endproof
\renewenvironment{proof}[1][\proofname]
  {\vspace{0.2cm}\begin{adjustwidth}{2em}{2em}
   \oldproof[#1]}
  {\endoldproof
\end{adjustwidth}}

% PROBLEM
\newenvironment{problem}[1]
  {\vspace{-0.1cm}\begin{coloredadjust}\begin{adjustwidth}{2em}{2em}
   \textit{Problem. }#1}
  {\endoldproof
\end{adjustwidth}\end{coloredadjust}\gspc}

%% DEFINITION
\newtcbtheorem[]{definition}{Definition}{enhanced,
  colback=blue!2, colframe=blue!50!black, coltitle=blue!70!black,
  fonttitle=\bfseries,
  boxrule=0pt, toprule=2pt, bottomrule=1pt, leftrule=1pt, rightrule=1pt, arc=1pt, outer arc=1pt,
  attach title to upper, title={#1},
}{dfn}

%% CRITERION
\newtcbtheorem[]{criterion}{Criterion}{enhanced,
  colback=magenta!3, colframe=magenta!60!black, coltitle=magenta!20!black,
  fonttitle=\bfseries,
  boxrule=4pt, toprule=2pt, bottomrule=1pt, leftrule=1pt, rightrule=1pt, arc=1pt, outer arc=1pt,
  attach title to upper, title={#1},
}{ctn}

%% COROLLARY
\newtcbtheorem[]{corollary}{Corollary}{enhanced,
  colback=orange!3, colframe=orange!80!black, coltitle=orange!50!black,
  fonttitle=\bfseries,
  boxrule=4pt, toprule=1pt, bottomrule=1pt, leftrule=1pt, rightrule=1pt, arc=5pt, outer arc=5pt,
  attach title to upper, title={#1},
}{ctn}

%% THEOREM
\newtcbtheorem[]{theorem}{Theorem}{enhanced,
  colback=green!2, colframe=green!40!black, coltitle=green!20!black,
  fonttitle=\bfseries,
  boxrule=4pt, toprule=3pt, bottomrule=1pt, leftrule=1pt, rightrule=1pt,arc=5pt, outer arc=5pt,
  attach title to upper, title={#1},
}{thm}

%%%% Page Setup %%%%%%%
\smallpage\sloppy
\pagestyle{fancy}
\lhead{\large{\textbf{HW1 - Theory of Computation}}}
\setlength{\headheight}{10pt}


% color commands
\newcommand{\cg}[1]{\color{OliveGreen}#1\color{white}}
\newcommand{\cb}[1]{\color{BlueViolet}#1\color{white}}


\newcommand{\cpr}[1]{\left(#1\right)}
%%%%%%%%%%%%% PHYSICS SPECIFIC COMMANDS

\newcommand{\vA}{\vec{\mathbf{A}}}
\newcommand{\vB}{\vec{\mathbf{B}}}
\newcommand{\vC}{\vec{\mathbf{C}}}
\newcommand{\vZ}{\vec{\mathbf{0}}}
\newcommand{\norm}{\vec{\mathbf{n}}}
\newcommand{\pvec}[1]{\left\langle#1\right\rangle}
\newcommand{\ar}{\vec{\mathbf{r}}}
\newcommand{\arhat}{\hat{\mathbf{r}}}
\newcommand{\vv}{\vec{\mathbf{v}}}
\newcommand{\delfull}{\pvec{\prt{}{x}, \prt{}{y},\prt{}{z}}}
\newcommand{\delinput}[3]{\pvec{\prt{}{x}\cpr{#1}, \prt{}{y}\cpr{#2},\prt{}{z}\cpr{#3}}}
\newcommand{\dps}{\displaystyle}
\newcommand{\delmatrix}[3]{\begin{bmatrix}\dps\hat x&\dps\hat y&\dps\hat z\\\dps\prt{}x&\dps\prt{}y&\dps\prt{}z\\\dps#1&\dps#2&\dps#3\end{bmatrix}}

%%%%%%%%%%%%%%%%%%%%%%%%%%%%
%%%%%%%%%%%%%%%%%%%%%%%%%%%%%%
%%%%%%%%%%%%%%%%%%%%%%%%%%%%%
%%%%%%%%%%%%%%%%%%%%%%%%%%%%%
%%%%%%%%%%%%%%%%%%%%%%%%%%%%%%
%%%%%%%%% begin doc


\begin{document}
\begin{enumerate}
  \item Let $A=\{1, 2,3 \}, B=\{1,2\}, C=\{2,4,5\}$
  \begin{enumerate}[itemsep=0pt, parsep=0pt]
    \item Is $A$ a subset of $B$? Explain.
    \begin{problem}
      No, it is not. Consider $3\in A$, $3\notin B$, and thus $A \nsubseteq B$. 
    \end{problem}
    \item Is $B$ a subset of $A$? Explain.
    \begin{problem}
      Yes, $B$ is a subset of $A$, because $\forall b\in B$, $b\in A$. 
    \end{problem}
    \item Is $B$ a subset of $C$?
    \begin{problem}
      No, because $1\in B$ and $1\notin C$. 
    \end{problem}
    \item What is $A\cup C$? Explain.
    \begin{problem}
     $A\cup C$ is defined as $A\cup C = \{x\mid x\in A\text{ or } x\in B\}$. Thus, the union of $A$ and $C$ is:
     \begin{align*}
      \{1, 2, 3, 4, 5\} \Rightarrow \text{ note must be distinct elements in a set}
     \end{align*}
    \end{problem}
    \item What is the power set of $A$? Explain.
    \begin{problem}
      The power set of $A$ is the set of all subsets of $A$. Thus, the power set of $A$, in order of cardinality, is:
      \begin{gather*}
        \{\{\},\\
        \{1\}, \{2\}, \{3\}\\
        \{1,2\}, \{1,3\}, \{2,3\}\\
        \{1,2,3\}\}
      \end{gather*}
    \end{problem}
  \end{enumerate}
  \item Consider the DFA pictured below
  \gimage{image.png}{6cm}
  \begin{enumerate}
    \item What is its alphabet?
    \begin{problem}
      It's alphabet is $\{0, 1\}$.
    \end{problem}
    \pagebreak
    \item What are its states? What is its start state? What are the accept states?
    \begin{problem}
      It's states are $\{q_1, q_2, q_3, q_4\}$.\gline{}
      It's start state is $q_4$.\gline{}
      The accept states are $\{q_3\}$. 
    \end{problem}
    \item Give five examples of strings represented by this DFA and five examples of strings now recognized by this DFA.
    \begin{problem}
      Five examples of strings represented by this DFA:
      \begin{enumerate}
        \item All binary strings of even length of the pattern $01\cdots$.
        \item I think that's the rule so I'm just going to give some examples:
        \begin{align*}
          \{01, 0101, 010101, 01010101, 0101010101\}
        \end{align*}
      \end{enumerate}
      Five examples of strings not recognized by this DFA:
      \begin{enumerate}
        \item All strings with consecutive ones. 
        \item All strings with consecutive zeros.
        \item All strings of odd length
        \item All strings that end in zero
        \item All strings that start with 1. 
      \end{enumerate}
      Some direct examples may look like:
      \begin{align*}
        \{11, 00, 01010101010, 010, 1\}
      \end{align*}
    \end{problem}
    \item Is the empty string recognized by this DFA?
    \begin{problem}
      No, it is not. This is because this string just remains in the start state, which is not an accepting state.
    \end{problem}
    \item Describe the language recognized by this DFA. Explain your answer.
    \begin{problem}
      I think I said this earlier, but I think the language is ``all strings that start with a zero, where every zero is followed by a one, where there are no consecutive 1s.'' In practice, this just looks like:
      \begin{align*}
        \{01, 0101, \ldots\}
      \end{align*}
      I might succinctly describe it as $A^*$ for $A = \{01\}$. I think this is the case because in the DFA diagram, whenever the binary string starts with 1, or there are two consecutive reads of the same number, you enter the infinite loop at $q_1$. Thus, it must follow an alternating pattern that starts with zero, and ends with 1 (because $q_2$ is not an accepting state).
    \end{problem}
  \end{enumerate}
  \pagebreak
  \item Consider the DFA pictured below. 
  \gimage{image2.png}{6cm}
  \begin{enumerate}
    \item What is its alphabet?
    \begin{problem}
      It's alphabet is $\{0, 1\}$. 
    \end{problem}
    \item What are its states? What is its start state? What are the accept states?
    \begin{problem}
      \begin{align*}
        &\text{Start State: } q_1\\
        &\text{States: }\{q_1, q_2, q_3, q_4, q_5\}\\
        &\text{Accept States: }\{q_3, q_5\}
      \end{align*}
    \end{problem}
    \item Give five examples of strings recognized by this DFA and five examples of strings not recognized by this DFA.
    \begin{problem}
      I think the pattern for this DFA is that a string only passes if:
      \begin{enumerate}
        \item The string ends in 1. 
        \item There are no consecutive 0s.
        \item No consecutive pair of 1s can be followed by a zero. 
      \end{enumerate}
      Thus, we get the example strings that \textit{satisfy} this DFA as being:
      \begin{align*}
        \{1, 01, 011, 0101, 01011\}
      \end{align*}
      And example strings that don't satisfy this DFA as:
      \begin{align*}
        \{0, 00, 10, 010, 001\}
      \end{align*}
    \end{problem}
    \item Is the empty string recognized by this DFA? Explain.
    \begin{problem}
      No, because the start state is not an accepting state.
    \end{problem}
    \item Describe the language recognized by this DFA
    \begin{problem}
      I think the language described by this DFA is parsing a binary string saying: ``This string must end in 1, contain no consecutive zeros, and ensure that all consecutive 1s are followed by another 1.''\gline{}
      Consider that $q_4$ is an infinite loop in a non-accepting state. If we look at the diagram, we can see that two zeros at any point will go to $q_4$ (the infinite loop). We also see that at any point, if we are at $q_5$, a zero will bring us to $q_4$. The only way to get to $q_5$ is if you come right from the start state, or are in a scenario where two ones have been parsed. Additionally, because $q_3$ and $q_5$ are our only accepting states (which the only way to get to is from a one), this tells us that the DFA is only accepting of strings that end in one, where no two consecutive ones are followed by a zero. \gline{}
      Note that we don't \textit{need} consecutive ones, because $\{1, 01\}$ are valid strings, but if they do appear, they cannot be followed by a zero. 
    \end{problem}
  \end{enumerate}
  \item Draw a DFA that recognizes the following languages:
  \begin{enumerate}
    \item The language of binary string with an odd number of characters
    \gimage{image3.png}{6cm}
    \item The language of binary strings that start and end with 1. 
    \gimage{image4.png}{6cm}
  \end{enumerate}
  \pagebreak
  \item Draw a DFA that recognizes the language of binary strings that start and end with 1 and have an even number of characters.
  \gimage{image5.png}{8cm}
\end{enumerate}
\end{document}